% Transplant-ready data section for the main paper.
% This is a fragment, not a full LaTeX document.

\section{Data and measurement}

This side evidence uses four public U.S. data objects: business counts by
employment size, business applications, CPS self-employment rates, and GEM
motive evidence. We do not collapse them into a single entrepreneurship
measure because they differ in unit of observation, denominator, and economic
margin. The main paper should treat the size-bin business counts and the
business-applications series as the primary motivation pair. The CPS
self-employment series and GEM evidence are best used as companion objects
that clarify measurement rather than as the central figures.

\subsection{Business counts by employment size}

Our stock-side series combines Census Nonemployer Statistics and Census
Business Dynamics Statistics to build annual national establishment counts from
1997 to 2022. The zero-employee series is total nonemployer establishments
(`NESTAB`). Positive-employment bins use BDS establishment counts (`ESTAB`) for
`1--4', `5--19', and `20+' employees, where `5--19' is formed by summing
`5--9' and `10--19', and `20+' is formed by summing `20--499' and `500+'.
We use establishment counts rather than firm counts because the zero-employee
object is naturally available as a nonemployer-establishment total, making
establishments the cleanest common unit across bins.

Recession timing is added from the FRED `USREC' indicator, aggregated to an
annual any-recession flag. In the Great Recession window from 2007 to 2010,
zero-employee establishments rise from 21,708,021 to 22,110,628 (1.85\%),
while `1--4' employee establishments fall 3.59\%, `5--19' employee
establishments fall 2.57\%, and `20+' employee establishments fall 5.82\%.
Employer establishments as a whole fall 3.58\%, while the sum of all tracked
bins rises only 0.54\%. This object should therefore be read as evidence on
business composition, not as a direct measure of startup flows or person-level
self-employment.

\subsection{Business applications}

Our flow-side companion uses Census Business Formation Statistics, accessed via
the seasonally adjusted FRED series for total business applications and
high-propensity business applications. Monthly data are available from 2005
through 2025 in the current pull and are aggregated to annual sums. We track
total applications and high-propensity applications separately and use the
high-propensity share as a simple measure of the composition of applicant
quality or employer orientation.

Over the 2007--2010 window, total applications fall from 2,652,392 to
2,497,405 (-5.84\%), while high-propensity applications fall from 1,493,428
to 1,163,075 (-22.12\%). The high-propensity share therefore declines from
56.30\% to 46.57\%. This flow object complements the size-bin stock evidence by
showing that business formation becomes weaker on the employer-oriented margin
even when zero-employee activity is relatively resilient.

\subsection{CPS self-employment companion}

To avoid treating nonemployer business counts as a person-level
self-employment series, we also build a CPS companion from BLS/FRED monthly
levels for unincorporated self-employed workers, incorporated self-employed
workers, and total employment among individuals age 16 and over. We compute
annual averages for 2000--2022 and report self-employment rates as shares of
total employment.

In the Great Recession window, the total self-employment rate falls from
11.06\% to 10.70\%, with both the unincorporated and incorporated components
also declining. This divergence is useful rather than problematic: it shows
that an establishment-side nonemployer count and a person-side
self-employment rate capture different margins and should not be combined
mechanically in a single headline statistic.

\subsection{GEM motive evidence}

Finally, we use the Global Entrepreneurship Monitor as a motive-based
companion rather than as a direct measure of the model object. The public U.S.
APS microdata currently yield a clean weighted TEA sample for 2019--2021, and
the official 2023--2024 U.S. GEM report provides a later public aggregate
update. Among weighted U.S. early-stage entrepreneurs in the 2019--2021
microdata, the share reporting that they started a business because jobs are
scarce is 40.9\% in 2019, 49.6\% in 2020, and 45.7\% in 2021. Over the same
years, the share citing wealth or very high income is 68.3\%, 64.7\%, and
73.4\%.

These shares are not mutually exclusive. That overlap is exactly why GEM is
useful here: it shows that necessity-type and opportunity-type signals can
coexist within the same observed entrepreneurial spell. We therefore use GEM
to motivate the paper's measurement problem, not to claim that survey motives
map one-for-one into the structural definition of necessity entrepreneurship.

Taken together, these objects imply two practical rules for the paper. First,
the main empirical motivation should center on the size-bin business-count
figure and the business-applications figure, because together they show a
shift toward low-scale activity and away from employer-oriented formation.
Second, CPS self-employment and GEM should remain companion evidence in prose,
footnotes, or appendix-style discussion. They sharpen the interpretation of
the main figures, but they do not collapse the paper's measurement problem
into a single observed statistic.
